\documentclass[11pt]{article}
% ===========================================================================
%  Shared preamble for all MPM2D worksheets — input right after \documentclass.
%  (Files beginning with "_" are skipped by mpm2d-worksheet-publish.mjs.)
% ===========================================================================
\usepackage[letterpaper,top=2.2cm,bottom=1.9cm,left=1.6cm,right=1.6cm,headheight=28pt,footskip=24pt]{geometry}
\usepackage{amsmath,amssymb}
\usepackage{cancel}
\usepackage[T1]{fontenc}
\usepackage[utf8]{inputenc}
\usepackage{lmodern}
\usepackage{xcolor}
\usepackage[most]{tcolorbox}
\usepackage{enumitem}
\usepackage{multicol}
\usepackage{fancyhdr}
\usepackage{tikz}
\usetikzlibrary{calc}
\usepackage{pgfplots}
\pgfplotsset{compat=1.18}

\definecolor{exblue}{HTML}{4A90E2}
\definecolor{exbg}{HTML}{E6F3FF}
\definecolor{qorange}{HTML}{E69138}
\definecolor{qbg}{HTML}{FFF7CC}
\definecolor{keygray}{HTML}{475569}
\definecolor{keybg}{HTML}{F1F5F9}
\definecolor{iagreen}{HTML}{14653B}
\definecolor{titleblue}{HTML}{1D4ED8}
% Rotating palette so each worked example gets its own colour.
\definecolor{ex0}{HTML}{2563A0}\definecolor{ex1}{HTML}{3B7D3B}\definecolor{ex2}{HTML}{8A5A00}
\definecolor{ex3}{HTML}{A3327A}\definecolor{ex4}{HTML}{6D28A3}\definecolor{ex5}{HTML}{B08900}
\definecolor{ex6}{HTML}{0E7490}\definecolor{ex7}{HTML}{9A3412}\definecolor{ex8}{HTML}{15803D}

\pagestyle{fancy}
\fancyhf{}
\fancyhead[L]{\footnotesize NAME: \underline{\hspace{3.4cm}}}
\fancyhead[C]{\textbf{Integration Academy}\\[-2pt]{\scriptsize Math Simplified \textperiodcentered{} Success Amplified}}
\fancyhead[R]{\footnotesize MPM2D}
\fancyfoot[L]{\scriptsize dr.merie@IntegrationAcademy.ca}
\fancyfoot[C]{\scriptsize\thepage}
\fancyfoot[R]{\scriptsize IntegrationAcademy.ca}
\renewcommand{\headrulewidth}{1.2pt}
\renewcommand{\footrulewidth}{0.4pt}
\setlength{\parindent}{0pt}
\setlength{\parskip}{4pt}

% Examples are NOT breakable, so each example + its solution stays on one page.
\newtcolorbox{example}[2]{colframe=#1,colback=#1!7,coltitle=white,colbacktitle=#1,fonttitle=\bfseries,title={#2},arc=3pt,boxrule=1pt}
\newtcolorbox{qbox}[1]{colframe=qorange,colback=qbg,coltitle=white,colbacktitle=qorange,fonttitle=\bfseries,title={#1},arc=3pt,breakable,boxrule=1pt}
\newtcolorbox{keybox}{colframe=keygray,colback=keybg,coltitle=white,colbacktitle=keygray,fonttitle=\bfseries,title={$\checkmark$\ Answer Key},arc=3pt,breakable,boxrule=1pt}
\newtcolorbox{recapbox}{colframe=keygray,colback=keybg,coltitle=white,colbacktitle=keygray,fonttitle=\bfseries,title={Question Recap},arc=3pt,breakable,boxrule=1pt}
\newtcolorbox{ideabox}{colframe=titleblue,colback=titleblue!6,coltitle=white,colbacktitle=titleblue,fonttitle=\bfseries,title={The Big Ideas},arc=3pt,breakable,boxrule=1.2pt}
\newtcolorbox{learnbox}{colframe=iagreen,colback=iagreen!5,coltitle=white,colbacktitle=iagreen,fonttitle=\bfseries,title={The Big Picture},arc=3pt,breakable,boxrule=1.2pt}
\newcommand{\lh}[1]{\par\smallskip{\bfseries\color{iagreen}#1}\par\nobreak\smallskip}

\newcommand{\soln}{\par\smallskip\textit{\textbf{Solution.}}\ }
\newcommand{\workspace}[1]{\par\vspace{3pt}\fcolorbox{black!30}{white}{\begin{minipage}[t][#1][t]{\dimexpr\linewidth-2\fboxsep\relax}\ \end{minipage}}\par}
\newcommand{\sech}[1]{\par\vspace{8pt}{\Large\bfseries\color{iagreen}#1}\par\nobreak\vspace{2pt}{\color{black!20}\hrule height1pt}\vspace{6pt}}

% A compact coordinate plot sized for an example box: \eplot{xmin}{xmax}{ymin}{ymax}{body}
\newcommand{\eplot}[5]{\begin{center}\begin{tikzpicture}
\begin{axis}[width=6.8cm,height=4.4cm,axis lines=middle,xlabel={\tiny$x$},ylabel={\tiny$y$},
grid=both,grid style={gray!16},xmin=#1,xmax=#2,ymin=#3,ymax=#4,
xtick distance=2,ytick distance=2,tick label style={font=\tiny},axis line style={-},samples=120]
#5
\end{axis}\end{tikzpicture}\end{center}}

% A sine-curve plot (degrees on x, 0..360): \splot{ymin}{ymax}{body}
\newcommand{\splot}[3]{\begin{center}\begin{tikzpicture}
\begin{axis}[width=8.6cm,height=4.6cm,axis lines=middle,xlabel={\tiny$x\,(^\circ)$},ylabel={\tiny$y$},
grid=both,grid style={gray!16},xmin=0,xmax=360,ymin=#1,ymax=#2,
xtick distance=90,ytick distance=1,tick label style={font=\tiny},axis line style={-},samples=180]
#3
\end{axis}\end{tikzpicture}\end{center}}

% A small coordinate plot: \plot{xmin}{xmax}{ymin}{ymax}{body}
\newcommand{\plot}[5]{\begin{center}\begin{tikzpicture}
\begin{axis}[width=8.4cm,height=6cm,axis lines=middle,xlabel={\scriptsize$x$},ylabel={\scriptsize$y$},
grid=both,grid style={gray!16},xmin=#1,xmax=#2,ymin=#3,ymax=#4,
xtick distance=2,ytick distance=2,tick label style={font=\tiny},axis line style={-}]
#5
\end{axis}\end{tikzpicture}\end{center}}

% Right triangle (angle theta at bottom-left, right angle at bottom-right):
%   \rtri{drawBase}{drawHeight}{oppLabel}{adjLabel}{hypLabel}
\newcommand{\rtri}[5]{\begin{center}\begin{tikzpicture}[scale=0.85]
\coordinate (A) at (0,0);\coordinate (B) at (#1,0);\coordinate (C) at (#1,#2);
\draw[thick,exblue] (A)--(B)--(C)--cycle;
\draw[black] ($(B)+(-0.32,0)$)--($(B)+(-0.32,0.32)$)--($(B)+(0,0.32)$);
\node[below] at ($(A)!0.5!(B)$){#4};
\node[right] at ($(B)!0.5!(C)$){#3};
\node[above left] at ($(A)!0.5!(C)$){#5};
\node at (0.7,0.22){$\theta$};
\end{tikzpicture}\end{center}}

% General (oblique) triangle with vertex labels and opposite-side labels:
%   \gtri{Alabel}{Blabel}{Clabel}{aSide}{bSide}{cSide}
\newcommand{\gtri}[6]{\begin{center}\begin{tikzpicture}[scale=0.85]
\coordinate (A) at (0,0);\coordinate (B) at (5,0);\coordinate (C) at (1.6,3);
\draw[thick,exblue] (A)--(B)--(C)--cycle;
\node[below left] at (A){#1};\node[below right] at (B){#2};\node[above] at (C){#3};
\node[below] at ($(A)!0.5!(B)$){#6};
\node[right] at ($(B)!0.5!(C)$){#4};
\node[left] at ($(A)!0.5!(C)$){#5};
\end{tikzpicture}\end{center}}

\fancyhead[R]{\footnotesize MCV4U \textperiodcentered{} 1.5}
\begin{document}

\begin{center}
{\LARGE\bfseries\color{titleblue}Worksheet 1.5 --- Derivative Using Definition}\\[2pt]
{\small Grade 12 Calculus and Vectors (MCV4U) \textperiodcentered{} Unit 1: Rates of Change and Limits}
\end{center}

\begin{ideabox}
$f'(x)=\displaystyle\lim_{h\to0}\dfrac{f(x+h)-f(x)}{h}$, with $x$ left general (unlike Lesson 1.4's single point). Expand, simplify to cancel $h$, then let $h\to0$.
\begin{itemize}[leftmargin=*,itemsep=2pt]
  \item $f'(x)=\displaystyle\lim_{h\to0}\dfrac{f(x+h)-f(x)}{h}$.
  \item Layout: Step 1 --- find $f(x+h)$ and expand it. Step 2 --- substitute into the definition, cross out the terms that cancel, factor out $h$, cancel it, and box the answer.
  \item Combine fractions / rationalize so $h$ cancels.
  \item Once $f'(x)$ is known, the slope at any point $a$ is just $f'(a)$ --- no limit needed.
\end{itemize}
\end{ideabox}

\begin{learnbox}
\lh{The definition}The \textbf{derivative} is $f'(x)=\displaystyle\lim_{h\to0}\dfrac{f(x+h)-f(x)}{h}$ — the same limit as Lesson 1.4, but with $x$ left general instead of one fixed point.
\lh{What it measures}$f'(x)$ gives the instantaneous rate of change, equivalently the slope of the tangent line, at \emph{every} $x$ at once.
\lh{Using the definition}Computing this limit directly is differentiating \emph{using the definition} (or \emph{first principles}), and it justifies the shortcut rules that follow.
\end{learnbox}

\sech{Worked Examples}

\begin{example}{ex0}{Example 1 --- Polynomial}
Differentiate $f(x)=x^2$ using the definition.\soln
\textbf{Step 1:} Find $f(x+h)$: $f(x+h)=(x+h)^2=x^2+2xh+h^2$.\\[3pt]
\textbf{Step 2:} Substitute into the definition, cross out what cancels, factor out $h$, and cancel it:
\begin{align*}
f'(x)&=\lim_{h\to0}\frac{f(x+h)-f(x)}{h}\\
&=\lim_{h\to0}\frac{(x^2+2xh+h^2)-(x)^2}{h}\\
&=\lim_{h\to0}\frac{\cancel{x^2}+2xh+h^2-\cancel{x^2}}{h}\\
&=\lim_{h\to0}\frac{\cancel{h}(2x+h)}{\cancel{h}}\\
&=\boxed{2x}
\end{align*}\eplot{-3}{3}{-6}{9}{\addplot[exblue,very thick,domain=-3:3]{x^2};\addplot[qorange,thick,domain=-3:3]{2*x};}
\end{example}

\begin{example}{ex1}{Example 2 --- Polynomial}
Differentiate $f(x)=3x^2-x$ using the definition.\soln
\textbf{Step 1:} Find $f(x+h)$: $f(x+h)=3(x+h)^2-(x+h)=3x^2+6xh+3h^2-x-h$.\\[3pt]
\textbf{Step 2:} Substitute into the definition, cross out what cancels, factor out $h$, and cancel it:
\begin{align*}
f'(x)&=\lim_{h\to0}\frac{f(x+h)-f(x)}{h}\\
&=\lim_{h\to0}\frac{(3x^2+6xh+3h^2-x-h)-(3x^2-x)}{h}\\
&=\lim_{h\to0}\frac{\cancel{3x^2}+6xh+3h^2-\cancel{x}-h-\cancel{3x^2}+\cancel{x}}{h}\\
&=\lim_{h\to0}\frac{\cancel{h}(6x+3h-1)}{\cancel{h}}\\
&=\boxed{6x-1}
\end{align*}
\end{example}

\begin{example}{ex2}{Example 3 --- Reciprocal}
Differentiate $f(x)=\dfrac1x$ using the definition.\soln
\textbf{Step 1:} Find $f(x+h)$: $f(x+h)=\dfrac{1}{x+h}$.\\[3pt]
\textbf{Step 2:} Substitute into the definition, combine the fractions over the common denominator $x(x+h)$, cross out what cancels, and cancel $h$:
\begin{align*}
f'(x)&=\lim_{h\to0}\frac{f(x+h)-f(x)}{h}\\
&=\lim_{h\to0}\dfrac{\dfrac{1}{x+h}-\dfrac{1}{x}}{h}\\
&=\lim_{h\to0}\dfrac{\dfrac{x-(x+h)}{x(x+h)}}{h}\\
&=\lim_{h\to0}\frac{\cancel{x}-\cancel{x}-h}{x(x+h)\cdot h}\\
&=\lim_{h\to0}\frac{-\cancel{h}}{x(x+h)\cancel{h}}\\
&=\lim_{h\to0}\frac{-1}{x(x+h)}\\
&=\boxed{-\frac{1}{x^2}}
\end{align*}\eplot{0.6}{3}{-3}{3}{\addplot[exblue,very thick,domain=0.6:3]{1/x};\addplot[qorange,thick,domain=0.6:3]{-1/x^2};}
\end{example}

\begin{example}{ex3}{Example 4 --- Using the formula}
From Example 1, $f'(x)=2x$. Find the tangent slope of $f(x)=x^2$ at $x=5$.\soln
\textbf{Step 1:} We already have $f'(x)=2x$ as a formula valid for every $x$ --- no need to redo the limit.\\[3pt]
\textbf{Step 2 (substitute):} $f'(5)=2(5)$.\\[3pt]
\textbf{Conclusion:} $10$ --- instant, once $f'(x)$ is known (compare to redoing the limit in Lesson 1.4).
\end{example}

\begin{example}{ex4}{Example 5 --- Radical}
Differentiate $f(x)=\sqrt{x}$ using the definition.\soln
\textbf{Step 1:} Find $f(x+h)$: $f(x+h)=\sqrt{x+h}$.\\[3pt]
\textbf{Step 2:} Substitute into the definition, multiply top and bottom by the conjugate $\sqrt{x+h}+\sqrt{x}$, cross out what cancels, and cancel $h$:
\begin{align*}
f'(x)&=\lim_{h\to0}\frac{f(x+h)-f(x)}{h}\\
&=\lim_{h\to0}\frac{\sqrt{x+h}-\sqrt{x}}{h}\\
&=\lim_{h\to0}\frac{\left(\sqrt{x+h}-\sqrt{x}\right)\left(\sqrt{x+h}+\sqrt{x}\right)}{h\left(\sqrt{x+h}+\sqrt{x}\right)}\\
&=\lim_{h\to0}\frac{\cancel{x}+h-\cancel{x}}{h\left(\sqrt{x+h}+\sqrt{x}\right)}\\
&=\lim_{h\to0}\frac{\cancel{h}}{\cancel{h}\left(\sqrt{x+h}+\sqrt{x}\right)}\\
&=\lim_{h\to0}\frac{1}{\sqrt{x+h}+\sqrt{x}}\\
&=\boxed{\frac{1}{2\sqrt{x}}}
\end{align*}\eplot{0}{4}{0}{2.5}{\addplot[exblue,very thick,domain=0:4]{sqrt(x)};\addplot[qorange,thick,domain=0.15:4]{1/(2*sqrt(x))};}
\end{example}

\begin{example}{ex5}{Example 6 --- Polynomial}
Differentiate $f(x)=x^2+5x$ using the definition.\soln
\textbf{Step 1:} Find $f(x+h)$: $f(x+h)=(x+h)^2+5(x+h)=x^2+2xh+h^2+5x+5h$.\\[3pt]
\textbf{Step 2:} Substitute into the definition, cross out what cancels, factor out $h$, and cancel it:
\begin{align*}
f'(x)&=\lim_{h\to0}\frac{f(x+h)-f(x)}{h}\\
&=\lim_{h\to0}\frac{(x^2+2xh+h^2+5x+5h)-(x^2+5x)}{h}\\
&=\lim_{h\to0}\frac{\cancel{x^2}+2xh+h^2+\cancel{5x}+5h-\cancel{x^2}-\cancel{5x}}{h}\\
&=\lim_{h\to0}\frac{\cancel{h}(2x+h+5)}{\cancel{h}}\\
&=\boxed{2x+5}
\end{align*}
\end{example}

\begin{example}{ex6}{Example 7 --- Polynomial}
Differentiate $f(x)=2x^2$ using the definition.\soln
\textbf{Step 1:} Find $f(x+h)$: $f(x+h)=2(x+h)^2=2x^2+4xh+2h^2$.\\[3pt]
\textbf{Step 2:} Substitute into the definition, cross out what cancels, factor out $h$, and cancel it:
\begin{align*}
f'(x)&=\lim_{h\to0}\frac{f(x+h)-f(x)}{h}\\
&=\lim_{h\to0}\frac{(2x^2+4xh+2h^2)-(2x^2)}{h}\\
&=\lim_{h\to0}\frac{\cancel{2x^2}+4xh+2h^2-\cancel{2x^2}}{h}\\
&=\lim_{h\to0}\frac{\cancel{h}(4x+2h)}{\cancel{h}}\\
&=\boxed{4x}
\end{align*}\eplot{-2}{2}{-8}{8}{\addplot[exblue,very thick,domain=-2:2]{2*x^2};\addplot[qorange,thick,domain=-2:2]{4*x};}
\end{example}

\begin{example}{ex7}{Example 8 --- Using the formula}
From Example 7, $f'(x)=4x$ for $f(x)=2x^2$. Find the tangent slope at $x=1$.\soln
\textbf{Step 1:} Substitute directly into the formula $f'(x)=4x$.\\[3pt]
\textbf{Step 2:} $f'(1)=4(1)$.\\[3pt]
\textbf{Conclusion:} $4$ --- no limit needed.
\end{example}

\begin{example}{ex8}{Example 9 --- Linear}
Differentiate $f(x)=5x-3$ using the definition.\soln
\textbf{Step 1:} Find $f(x+h)$: $f(x+h)=5(x+h)-3=5x+5h-3$.\\[3pt]
\textbf{Step 2:} Substitute into the definition, cross out what cancels, and cancel $h$:
\begin{align*}
f'(x)&=\lim_{h\to0}\frac{f(x+h)-f(x)}{h}\\
&=\lim_{h\to0}\frac{(5x+5h-3)-(5x-3)}{h}\\
&=\lim_{h\to0}\frac{\cancel{5x}+5h-\cancel{3}-\cancel{5x}+\cancel{3}}{h}\\
&=\lim_{h\to0}\frac{5\cancel{h}}{\cancel{h}}\\
&=\boxed{5}
\end{align*}
A line's derivative is its constant slope.
\end{example}

\clearpage
\sech{Your Turn --- Show Your Work}
For each question, show all steps clearly in the space provided.

\begin{qbox}{Question 1}
Using the definition, differentiate $f(x)=x^2$.\workspace{2.4cm}
\end{qbox}

\begin{qbox}{Question 2}
Using the definition, differentiate $f(x)=x^2+5x$.\workspace{2.4cm}
\end{qbox}

\begin{qbox}{Question 3}
Using the definition, differentiate $f(x)=2x^2$.\workspace{2.4cm}
\end{qbox}

\begin{qbox}{Question 4}
Given $f'(x)=2x$ for $f(x)=x^2$, find the tangent slope at $x=1$.\workspace{2.4cm}
\end{qbox}

\begin{qbox}{Question 5}
Using the definition, differentiate $f(x)=5x-3$.\workspace{2.4cm}
\end{qbox}

\begin{qbox}{Question 6}
Using the definition, differentiate $f(x)=4x^2-x$.\workspace{2.4cm}
\end{qbox}

\begin{qbox}{Question 7}
Using the definition, differentiate $f(x)=\dfrac1x$.\workspace{2.4cm}
\end{qbox}

\begin{qbox}{Question 8}
Given $f'(x)=2x$ for $f(x)=x^2$, find the tangent slope at $x=2$.\workspace{2.4cm}
\end{qbox}

\begin{qbox}{Question 9}
Using the definition, differentiate $f(x)=x^2-3x$.\workspace{2.4cm}
\end{qbox}

\begin{qbox}{Question 10}
Using the definition, differentiate $f(x)=7$.\workspace{2.4cm}
\end{qbox}

\begin{qbox}{Question 11}
Given $f'(x)=2x$ for $f(x)=x^2$, find the tangent slope at $x=-1$.\workspace{2.4cm}
\end{qbox}

\begin{qbox}{Question 12}
Using the definition, differentiate $f(x)=3x^2$.\workspace{2.4cm}
\end{qbox}

\begin{qbox}{Question 13 --- Challenge}
Using the definition, differentiate $f(x)=\sqrt{x}$ (rationalize).\workspace{3.5cm}
\end{qbox}

\clearpage
\sech{Question Recap \& Answer Key}

\begin{recapbox}
\begin{multicols}{2}
\small
\begin{enumerate}[leftmargin=*,itemsep=3pt]
  \item Using the definition, differentiate $f(x)=x^2$.
  \item Using the definition, differentiate $f(x)=x^2+5x$.
  \item Using the definition, differentiate $f(x)=2x^2$.
  \item Given $f'(x)=2x$ for $f(x)=x^2$, find the tangent slope at $x=1$.
  \item Using the definition, differentiate $f(x)=5x-3$.
  \item Using the definition, differentiate $f(x)=4x^2-x$.
  \item Using the definition, differentiate $f(x)=\dfrac1x$.
  \item Given $f'(x)=2x$ for $f(x)=x^2$, find the tangent slope at $x=2$.
  \item Using the definition, differentiate $f(x)=x^2-3x$.
  \item Using the definition, differentiate $f(x)=7$.
  \item Given $f'(x)=2x$ for $f(x)=x^2$, find the tangent slope at $x=-1$.
  \item Using the definition, differentiate $f(x)=3x^2$.
  \item Using the definition, differentiate $f(x)=\sqrt{x}$ (rationalize).
\end{enumerate}
\end{multicols}
\end{recapbox}

\begin{keybox}
\begin{multicols}{2}
\small
\begin{enumerate}[leftmargin=*,itemsep=3pt]
  \item $2x$
  \item $2x+5$
  \item $4x$
  \item $2$
  \item $5$
  \item $8x-1$
  \item $-\tfrac1{x^2}$
  \item $4$
  \item $2x-3$
  \item $0$
  \item $-2$
  \item $6x$
  \item $\tfrac1{2\sqrt{x}}$
\end{enumerate}
\end{multicols}
\end{keybox}

\end{document}
